\documentclass[compress, 8pt, aspectratio=169]{beamer}

\usepackage{presentationtemplate}
\usepackage[askip=3mm, bskip=3mm]{terminal}
\usepackage[linenosfontsize=\tiny, askip=3mm, bskip=3mm]{mylisting}
\usepackage{tabularray}

% В окружении сборки старый pygments:
% символ `|` из count.cpp (режим in|out) он
% классифицирует как ошибочный токен,
% который minted по умолчанию рендерит
% красной рамкой. Эта опция заставляет
% такие токены рендериться
% обычным моноширинным текстом.
\setminted{ignorelexererrors=true}

\tcbset{
    listingbase17/.style={
        listing engine=minted,
        minted options={
            fontsize=\footnotesize,
            breaklines,
            breakanywhere,
            linenos,
            escapeinside=||,
            numbersep=2pt
        },
        listing only,
        left skip=0mm,
        right skip=0mm,
        after skip=3mm,
        before skip=3mm,
        left=0mm,
        right=0mm,
        top=0mm,
        bottom=0mm,
        boxrule=0.02cm,
        colback=white,
        colframe=black,
        colbacktitle=gray!10!white,
        coltitle=black,
        sharp corners
    }
}

\renewtcbinputlisting{\myinputlisting}[3][]{%
    listing file={#2#3},
    title=#3,
    fonttitle=\it\footnotesize,
    listingbase17,
    #1
}

\NewTCBListing{terminalwindow17}{ O{} }{
    listing engine=minted,
    title=>\_\null\hfill x,
    minted options={
        fontsize=\footnotesize,
        breaklines,
        breakanywhere,
        escapeinside=!!,
        breaksymbolleft=
    },
    minted language=text,
    listing only,
    left skip=0mm,
    right skip=0mm,
    after skip=3mm,
    before skip=3mm,
    left=0mm,
    right=0mm,
    top=0mm,
    bottom=0mm,
    boxrule=0.02cm,
    colback=white,
    colframe=black,
    colbacktitle=black,
    fonttitle=\ttfamily\footnotesize,
    #1
}

\title{Стандартная библиотека: потоки}

\begin{document}

  \frame[plain]{\titlepage}

  \begin{frame}

    \frametitle{Потоки стандартной библиотеки}

    Поток представляет собой последовательность символов:
    одни и те же операторы \texttt{>>} (чтение)
    и \texttt{<<} (запись) работают с любым потоком.

    \vspace{\baselineskip}
    \begin{tblr}{
      colspec = {|X[0.22,l]|X[0.58,l]|X[0.20,l]|},
      row{1} = {font=\bfseries},
      hlines, vlines,
      rowsep = 0.8ex,
      colsep = 1.0ex,
    }
      Заголовок & Потоки & Назначение \\
      \texttt{<iostream>} & \texttt{std::cin}, \texttt{std::cout}, \texttt{std::cerr} & консоль \\
      \texttt{<sstream>} & \texttt{std::istringstream}, \texttt{std::ostringstream}, \texttt{std::stringstream} & строка \\
      \texttt{<fstream>} & \texttt{std::ifstream}, \texttt{std::ofstream}, \texttt{std::fstream} & файл \\
    \end{tblr}

  \end{frame}

  \begin{frame}[fragile]

    \frametitle{\texttt{std::cin} и \texttt{std::cout}}

    \begin{itemize}

      \item Стандартные объекты \texttt{std::cin} и \texttt{std::cout}
        объявлены в заголовке \texttt{<iostream>}.

      \item \texttt{>>} читает значения (пропуская пробельные),
        \texttt{<<} записывает.

    \end{itemize}

    \myinputlisting[minted language=cpp]
        {Presentations/17-streams/}
        {echo.cpp}

    \begin{terminalwindow17}
!\shellcommand{g++ -std=c++23 echo.cpp -o echo}!
!\shellcommand{printf 'Ada 36\textbackslash n' | ./echo}!
Ada, 36
    \end{terminalwindow17}

  \end{frame}

  \begin{frame}[fragile]

    \frametitle{Манипуляторы}

    \begin{itemize}

      \item Манипуляторы из \texttt{<iomanip>} меняют формат вывода.

      \item \texttt{std::setw} задаёт ширину поля, \texttt{std::fixed}
        включает формат с фиксированной десятичной точкой,
        \texttt{std::setprecision} задаёт число знаков после точки.

    \end{itemize}

    \myinputlisting[minted language=cpp]
        {Presentations/17-streams/}
        {manip.cpp}

    \begin{terminalwindow17}
!\shellcommand{g++ -std=c++23 manip.cpp -o manip \&\& ./manip}!
  1: 3.14
    \end{terminalwindow17}

  \end{frame}

  \begin{frame}[fragile]

    \frametitle{\texttt{operator<<} для своего типа}

    \myinputlisting[minted language=cpp, title=point\_out.cpp]
        {Presentations/17-streams/}
        {point_out.cpp}

    \begin{terminalwindow17}
!\shellcommand{g++ -std=c++23 point\_out.cpp -o point\_out \&\& ./point\_out}!
(1, 2) -> (3, 4)
    \end{terminalwindow17}

  \end{frame}

  \begin{frame}[fragile]

    \frametitle{\texttt{operator>>} для своего типа}

    \begin{itemize}

      \item Симметрично: \texttt{Point\&} (без \texttt{const}),
        поскольку объект заполняется; возвращаемый
        \texttt{in} поддерживает цепочки.

      \item Состояние потока передаётся: работает идиома
        \texttt{while (in >> p)} (слайд 11).

    \end{itemize}

    \myinputlisting[minted language=cpp, title=point\_in.cpp]
        {Presentations/17-streams/}
        {point_in.cpp}

    \begin{terminalwindow17}
!\shellcommand{g++ -std=c++23 point\_in.cpp -o point\_in}!
!\shellcommand{printf '3 4\textbackslash n' | ./point\_in}!
3 4
    \end{terminalwindow17}

  \end{frame}

  \begin{frame}[fragile]

    \frametitle{\texttt{ostringstream}: запись в строку}

    \begin{itemize}

      \item \texttt{std::ostringstream} из заголовка
        \texttt{<sstream>} пишет в строку.

      \item Оператор \texttt{<<} тот же; результат читается
        методом \texttt{str()}.

    \end{itemize}

    \myinputlisting[minted language=cpp]
        {Presentations/17-streams/}
        {ostr.cpp}

    \begin{terminalwindow17}
!\shellcommand{g++ -std=c++23 ostr.cpp -o ostr \&\& ./ostr}!
answer: 42
    \end{terminalwindow17}

  \end{frame}

  \begin{frame}[fragile]

    \frametitle{\texttt{istringstream}: чтение из строки}

    \begin{itemize}

      \item \texttt{std::istringstream} из заголовка
        \texttt{<sstream>} читает из строки.

      \item \texttt{>>} разбирает значения разных типов.

    \end{itemize}

    \myinputlisting[minted language=cpp]
        {Presentations/17-streams/}
        {istr.cpp}

    \begin{terminalwindow17}
!\shellcommand{g++ -std=c++23 istr.cpp -o istr \&\& ./istr}!
42 hello
    \end{terminalwindow17}

  \end{frame}

  \begin{frame}[fragile]

    \frametitle{\texttt{stringstream}: чтение и запись}

    \begin{itemize}

      \item \texttt{std::stringstream} сочетает \texttt{>>}
        и \texttt{<<} в одном объекте.

      \item После чтения до конца (eofbit) запись невозможна:
        требуется \texttt{clear()}.

    \end{itemize}

    \myinputlisting[minted language=cpp]
        {Presentations/17-streams/}
        {ss.cpp}

    \begin{terminalwindow17}
!\shellcommand{g++ -std=c++23 ss.cpp -o ss \&\& ./ss}!
sum: 30
    \end{terminalwindow17}

  \end{frame}

  \begin{frame}[fragile]

    \frametitle{\texttt{ofstream}: запись в файл}

    \begin{itemize}

      \item \texttt{std::ofstream} из заголовка
        \texttt{<fstream>} пишет в файл.

      \item По умолчанию файл создаётся заново (существующее
        содержимое удаляется); оператор \texttt{<<} тот же,
        что у \texttt{std::cout}.

    \end{itemize}

    \myinputlisting[minted language=cpp, title=write\_file.cpp]
        {Presentations/17-streams/}
        {write_file.cpp}

    \begin{terminalwindow17}
!\shellcommand{g++ -std=c++23 write\_file.cpp -o write \&\& ./write}!
!\shellcommand{cat data.txt}!
1 2 3
    \end{terminalwindow17}

  \end{frame}

  \begin{frame}[fragile]

    \frametitle{\texttt{ifstream}: чтение из файла}

    \begin{itemize}

      \item \texttt{std::ifstream} из заголовка
        \texttt{<fstream>} читает из файла.

      \item Идиома \texttt{while (in >> n)}: цикл завершается
        при неудачном чтении (конец файла).

    \end{itemize}

    \myinputlisting[minted language=cpp, title=read\_file.cpp]
        {Presentations/17-streams/}
        {read_file.cpp}

    \begin{terminalwindow17}
!\shellcommand{g++ -std=c++23 read\_file.cpp -o read \&\& ./read}!
1
2
3
    \end{terminalwindow17}

  \end{frame}

  \begin{frame}[fragile]

    \frametitle{\texttt{fstream}: чтение и запись}

    \myinputlisting[minted language=cpp]
        {Presentations/17-streams/}
        {count.cpp}

    \begin{terminalwindow17}
!\shellcommand{printf 3 > count.txt}!
!\shellcommand{g++ -std=c++23 count.cpp -o count \&\& ./count}!
4
!\shellcommand{cat count.txt}!
3 4
!\shellcommand{./count \&\& cat count.txt}!
4
3 4 4
    \end{terminalwindow17}

  \end{frame}

  \begin{frame}

    \frametitle{Режимы открытия файла}

    Состояние потока: \texttt{good()}, \texttt{fail()}.

    \vspace{\baselineskip}
    \begin{tblr}{
      colspec = {|X[0.30,l]|X[0.30,l]|X[0.40,l]|},
      row{1} = {font=\bfseries},
      hlines, vlines,
      rowsep = 0.8ex,
      colsep = 1.0ex,
    }
      Класс & Режим по умолчанию & Файла нет \\
      \texttt{std::ifstream} & \texttt{in} & устанавливается \texttt{failbit} \\
      \texttt{std::ofstream} & \texttt{out} & файл создаётся, существующее содержимое удаляется \\
      \texttt{std::fstream} & \texttt{in | out} & устанавливается \texttt{failbit} \\
    \end{tblr}

  \end{frame}

  \begin{frame}[fragile]

    \frametitle{\texttt{ifstream} или \texttt{ofstream}?}

    \myinputlisting[minted language=cpp, title=open\_check.cpp]
        {Presentations/17-streams/}
        {open_check.cpp}

    \begin{terminalwindow17}
!\shellcommand{g++ -std=c++23 open\_check.cpp -o check \&\& ./check}!
ofstream: true
ifstream: false
!\shellcommand{ls new.txt}!
new.txt
    \end{terminalwindow17}

  \end{frame}

\end{document}
